\subsection*{二分之一从一块三角形出现}
设 $F'=f$，想估计 $F(a+h)-F(a)$。第一步把高度固定为 $f(a)$，得到矩形面积 $f(a)h$。但若高度随着输入增加，它的变化也会贡献面积。把高度函数用直线近似，末端比起点高约 $f'(a)h$，新增区域近似三角形。

三角形底为 $h$、高为 $f'(a)h$，所以修正是 $\frac12 f'(a)h^2$。二阶项的二分之一因此可以看成变化高度的平均效果。若斜率为负，修正也带负号；有向面积让同一公式继续成立。

\subsection*{一次累积把幂次提高一级}
高度函数的一阶变化形成累积函数的二阶修正。若继续近似高度本身的弯曲，再积分，就会产生三阶项。积分每次提高幂次并除以新的次数，这与泰勒系数里的阶乘相呼应。

对 $F(x)=x^2$，导数 $f(x)=2x$ 本来就是直线，所以矩形加三角形已经精确，没有更高误差。对一般函数，三角形只模拟局部高度变化，还需要估计曲线与这条直线之间的面积差。

把本节与上一节对照：一种办法逐阶匹配导数，另一种从累积几何解释修正。两种路线得到同一局部表达，但适用的光滑条件仍要保留。看图时不要把“画起来像三角形”直接升级为精确等式，要先判断高度函数是否真的线性。
